\documentclass[11pt]{article}

\usepackage[T1]{fontenc}
\usepackage{amsmath,amssymb,amsthm}
\usepackage{booktabs}
\usepackage{geometry}
\usepackage{hyperref}
\usepackage{microtype}
\usepackage{xcolor}

\geometry{margin=0.93in}
\hypersetup{
  colorlinks=true,
  linkcolor=blue!50!black,
  citecolor=blue!50!black,
  urlcolor=blue!50!black,
  pdftitle={A Counterexample to Minimum Fano-Projection Selection},
  pdfauthor={Atharva Vaidya}
}

\newtheorem{theorem}{Theorem}[section]
\newtheorem{lemma}[theorem]{Lemma}
\newtheorem{proposition}[theorem]{Proposition}
\newtheorem{corollary}[theorem]{Corollary}
\theoremstyle{definition}
\newtheorem{definition}[theorem]{Definition}
\theoremstyle{remark}
\newtheorem{remark}[theorem]{Remark}

\newcommand{\F}{\mathbb F}
\newcommand{\supp}{\operatorname{supp}}
\newcommand{\symdiff}{\mathbin{\triangle}}
\newcommand{\FiveCDC}{\textup{FiveCDC}}
\newcommand{\bd}{\partial}
\newcommand{\dotcupx}{\mathbin{\dot\cup}}

\title{A Counterexample to Minimum Fano-Projection Selection\\
\large for Five-Cycle Double Covers}
\author{Atharva Vaidya\\
\small AI-assisted research note for independent human verification}
\date{July 29, 2026}

\begin{document}
\maketitle

\begin{center}
\fcolorbox{red!65!black}{red!4}{
\parbox{0.91\textwidth}{
\textbf{Scope and resolution status.}
This paper disproves a proposed \emph{proof method}: selecting a
cardinality-minimum binary projection of a nowhere-zero
\(\F_2^3\)-flow need not yield the component parity required by the
Hu\v{s}ek--\v{S}\'amal criterion.  The constructed graph itself has an
explicit five-cycle double cover.  No counterexample to \FiveCDC, and no
resolution of \FiveCDC, is claimed.
}}
\end{center}

\begin{abstract}
Write a nowhere-zero \(\F_2^3\)-flow on a cubic graph as
\(f=(h,s)\), where \(h\) is a binary flow and
\(s\) is an \(\F_2^2\)-flow.  Call \(h\) an extendable projection.
A recent flow criterion for the five-cycle double cover conjecture
requires an extendable projection with a component-parity property,
called clean here.  A natural proposed selection principle was that some
cardinality-minimum extendable projection should be cleanable.

We give a finite simple connected bridgeless cubic counterexample.
The graph has \(162\) vertices and \(243\) edges.  Its unique globally
minimum extendable projection has \(54\) edges, is the disjoint union of
two \(27\)-circuits, and has no clean extension.  The construction starts
from an \(18\)-vertex uncleanable projection and replaces eight support
edges by an \(18\)-vertex two-pole.  Exact cycle-space enumeration proves
that this two-pole has effective costs \(w_0=7\) and \(w_1=6\), with a
unique minimum support in the \(1\)-terminal state.  A transparent
weighted contraction argument then proves global minimality and
uniqueness; uncleanability descends under contraction.

The negative result is strictly about the selection method.  We give a
human compositional five-cycle-double-cover construction for the same
graph and a second literal assignment for the mandatory SAT/XOR
encoding.  Two independently written standard-library checkers reproduce
the conditional minima, all \(2^{14}\) lock placements, the
uncleanability census, the canonical graph, and both positive cover
certificates.
\end{abstract}

\section{Conventions and the selection principle}

All graphs in the theorem are finite, loopless, and cubic.  Parallel
edges would cause no difficulty in the flow statements, but the witness
is simple.  Edges are distinct objects.  A \emph{binary cycle} is an
Eulerian edge set, possibly empty or disconnected.  Thus ``five
Eulerian subgraphs'' and ``at most five, after appending empty members''
are equivalent conventions.

Over a group of characteristic two, orientations may be suppressed:
a flow \(f:E(G)\to\F_2^k\) satisfies
\[
                    \sum_{e\ni v} f(e)=0
                    \qquad(v\in V(G)).
\]
It is nowhere zero if \(f(e)\ne0\) for every edge.  Put
\(K=\F_2^2\) and write
\[
                         f=(h,s):E(G)\longrightarrow
                         \F_2\times K.                    \tag{1}
\]
The support of the binary flow \(h\) is again denoted by \(h\).

\begin{definition}
A binary cycle \(h\) is \emph{extendable} if there is a \(K\)-flow
\(s\) for which \((h,s)\) is nowhere zero.  Its minimum possible
cardinality is
\[
        \mu(G)=\min\{|h|:h\text{ is an extendable projection}\}.
\]
\end{definition}

Fix an extension \((h,s)\).  For \(c\in K\), let
\[
                         M_c=\{e\in h:s(e)=c\}.            \tag{2}
\]
At a cubic vertex no two incident edges can have the same nonzero full
value, so every \(M_c\) is a matching.  The four matchings partition
\(h\).

Let \(W\) be a component of \(G-h\), and put
\[
       n_c(W)=|\delta(W)\cap M_c|\pmod2.
\]
Cut conservation for the three coordinates implies that the four
parities \(n_c(W)\) are equal.  This gives the following terminology.

\begin{definition}
An extension \((h,s)\) is \emph{clean} if \(n_c(W)=0\) for every
component \(W\) of \(G-h\), for one and hence every \(c\in K\).
The projection \(h\) is \emph{cleanable} if it has a clean extension.
\end{definition}

Hu\v{s}ek and \v{S}\'amal prove that, for cubic graphs, existence of a
clean extension is equivalent to existence of a five-cycle double
cover~\cite[Theorem~3.16 and Conjecture~3.19]{HusekSamal2026}.
Their formulation is related to Hoffmann-Ostenhof's necessary and
sufficient condition for a prescribed \(2\)-regular subgraph to occur
in a five-cycle double cover~\cite{HoffmannOstenhof2013}.
The selection statement refuted here is
\begin{equation}
\begin{split}
\textit{Every bridgeless cubic graph has a cleanable
cardinality-minimum}\\[-2mm]
\textit{extendable projection.}
\end{split}
\tag{S}
\end{equation}
Statement (S) is not the \FiveCDC\ conjecture itself.

\section{The full-flow master lemma}

The following elementary characterization makes ``globally minimum''
literal; in particular, it compares different low flows, not merely
four affine rebases of one fixed extension.

\begin{lemma}[full-flow master lemma]\label{lem:master}
Let \(G\) be a finite loopless cubic graph.  For a \(K\)-flow \(s'\),
write \(Z(s')=\{e:s'(e)=0\}\).  Then
\[
 \mu(G)=
 \min_{\substack{s'\text{ a }K\text{-flow}\\
                  J\subseteq E-Z(s')\\
                  \bd J=\bd Z(s')}}
       \bigl(|Z(s')|+|J|\bigr),                            \tag{3}
\]
where an infeasible inner problem has value \(+\infty\).
Every feasible zero set \(Z(s')\) is a matching.
\end{lemma}

\begin{proof}
Let \(M=Z(s')\), and suppose
\(J\subseteq E-M\) with \(\bd J=\bd M\).  Then
\(h'=M\dotcupx J\) has even degree at every vertex.  The flow
\((h',s')\) is nowhere zero: the low value vanishes exactly on \(M\),
where the first coordinate is one.

Conversely, from any nowhere-zero \((h',s')\), its exact low zero set
\(M\) is contained in \(h'\).  Put \(J=h'-M\).  The sets are disjoint
and, since \(h'\) is binary, \(\bd J=\bd M\).  These constructions are
inverse and preserve the objective \(|h'|=|M|+|J|\), proving (3).

At a cubic vertex the zero-degree of a \(K\)-flow is \(0,1\), or \(3\):
two zero incident values force the third to be zero.  Degree \(3\) is
incompatible with a disjoint \(J\) having the same boundary.  Hence a
feasible \(M\) is a matching.
\end{proof}

For fixed \(M\), the inner problem is the ordinary shortest
\(\bd M\)-join problem in \(G-M\), with its standard odd-cut dual.  The
same dictionary isolates cleanliness.

\begin{lemma}[the missing second join]\label{lem:cleanjoin}
For a feasible triple \((s',M,J_1)\) in Lemma~\ref{lem:master}, put
\(h=M\dotcupx J_1\).  The extension \((h,s')\) is clean if and only if
there is a second \(\bd M\)-join
\[
                          J_2\subseteq G-h.                \tag{4}
\]
\end{lemma}

\begin{proof}
A \(\bd M\)-join exists in a graph precisely when each component
contains an even number of vertices of \(\bd M\).  For a component
\(W\) of \(G-h\),
\[
 |W\cap\bd M|\equiv|M\cap\delta(W)|=n_0(W)\pmod2.
\]
The four cut-colour parities are equal by conservation, so their common
vanishing is exactly the component condition for the join in (4).
\end{proof}

Thus (S) asserts that some optimum branch of the disjunctive
optimization (3) admits the extra packing (4).  The graph below shows
that this need not happen.

\section{An uncleanable base projection}

Let \(B\) be the \(18\)-vertex cubic graph whose first fourteen edges
are two disjoint \(7\)-circuits:
\[
 (i,i+1\bmod7)\ (0\le i<7),\qquad
 (7+i,7+(i+1\bmod7))\ (0\le i<7).
\]
The remaining thirteen edges, in order, are
\[
\begin{gathered}
(0,11),(1,9),(2,12),(3,10),\\
(14,4),(14,5),(14,15),(15,6),(15,16),\\
(16,13),(16,17),(17,7),(17,8).
\end{gathered}                                             \tag{5}
\]
Let \(T=\{0,\ldots,13\}\), the union of the two \(7\)-circuits.
The low word
\[
             \texttt{3231320|1320320|3112213212312}       \tag{6}
\]
extends \(T\) to a nowhere-zero \(\F_2^3\)-flow.

\begin{proposition}[base obstruction]\label{prop:base}
The projection \(T\) has \(15{,}360\) ordered low-coordinate extensions
and none is clean.
\end{proposition}

\begin{proof}[Exact finite proof]
The cycle-space dimension of \(B\) is
\(27-18+1=10\).  Enumerate its \(1{,}024\) binary cycles \(p,q\).
The pair extends \(T\) precisely when
\[
                         E(B)-T\subseteq p\cup q.          \tag{7}
\]
There are \(15{,}360\) ordered pairs satisfying (7).  The five
components of \(B-T\), restricted to vertices \(0,\ldots,13\), have
label word
\[
                         \texttt{01234444413024}.          \tag{8}
\]
For every retained pair and every component, the checker xors the four
cut-colour incidence rows.  No pair makes all rows zero.  This is a
complete enumeration of all extensions by (7), not a sample or a
solver heuristic.
\end{proof}

The projection \(T\) is not globally minimum on \(B\).  We next make it
uniquely minimum without changing the contraction obstruction.

\section{A strict two-pole lock}

Let \(S\) be the simple bridgeless cubic graph with graph6 encoding
\[
             \texttt{Q???C@?K?WEOM?\_aGo?I\_\_W@G\_?}.    \tag{9}
\]
Its ordered edges are
\[
\begin{split}
&(0,7),(1,8),(2,9),(3,9),(4,10),(5,10),(2,11),(3,11),(6,11),\\
&(2,12),(3,12),(4,12),(0,13),(6,13),(10,13),(1,14),(5,14),\\
&(6,14),(5,15),(7,15),(9,15),(0,16),(7,16),(8,16),\\
&(1,17),(4,17),(8,17).
\end{split}                                                \tag{10}
\]
Distinguish edge \(e_*=e_5=(5,10)\).  For \(b\in\{0,1\}\), let
\(a_b\) be the minimum first-coordinate support size of a nowhere-zero
\(\F_2^3\)-flow on \(S\), conditional on the first bit of \(e_*\)
being \(b\).

\begin{lemma}[strict lock]\label{lem:lock}
\[
                            a_0=7,\qquad a_1=5.            \tag{11}
\]
The \(a_1\)-minimum support is unique and equals
\[
                           \{5,13,14,16,17\},              \tag{12}
\]
the \(5\)-circuit \(5-10-13-6-14-5\).
\end{lemma}

\begin{proof}[Exact finite proof]
The cycle space of \(S\) again has dimension ten.  A binary cycle \(h\)
is extendable precisely when two cycles \(p,q\) satisfy
\(h\cup p\cup q=E(S)\).  Enumerating the \(1{,}024\) choices of \(h\)
and the \(70{,}780\) distinct exact zero sets generated by ordered
\((p,q)\) gives (11).  There are nine \(a_0\)-minimum supports, with
\(6{,}336\) ordered extensions.  In the \(b=1\) branch, (12) is the
only minimum support and has \(240\) ordered extensions.  A second
checker constructs the cycle space by incidence-matrix elimination,
rather than fundamental circuits, and obtains the same counts.
\end{proof}

Delete \(e_*\), expose its endpoints, and attach one new terminal link
at each endpoint when the pole replaces an edge.  If the common
terminal first bit is \(0\), the least projected cost is
\[
                              w_0=a_0=7.
\]
If it is \(1\), deletion removes the marked projected edge and the two
links add two, so
\[
                              w_1=a_1-1+2=6.               \tag{13}
\]
This is the strict parity lock \(w_0>w_1\).

For completeness, the frozen full flow
\[
\begin{split}
(4,6,2,6,4,5,4,2,6,6,4,2,2,3,1,2,7,5,2,6,4,6,2,4,4,6,2)
\end{split}                                                \tag{14}
\]
attains (12), with the first coordinate represented by the least
significant bit.  Its marked value is \((1,2)\in\F_2\times K\).
For any desired odd terminal value \((1,t)\), the invertible map
\[
                 \Phi_t(x,\ell)=(x,\ell+x(2+t))           \tag{15}
\]
preserves the first coordinate and sends \((1,2)\) to \((1,t)\).
Thus the same unique projected pole support lifts every odd terminal
value required by the base flow.

\section{Eight locks and the 162-vertex theorem}

For \(L\subseteq T\), replace the base edges in \(L\) by disjoint fresh
copies of the pole \(S-e_*\); leave every other base edge unchanged.
The target terminal support \(T\) has cost
\[
                              14+5|L|.                    \tag{16}
\]
For another extendable terminal projection \(H\), put
\(O=T-H\) and \(A=H-T\).  The cost difference is
\[
 \Delta_L(H)
   =|O\cap L|-|O-L|+|A|
   =2|O\cap L|-|O|+|A|.                                  \tag{17}
\]
Indeed, omitting a locked target edge costs one, omitting an unlocked
target edge saves one, and adding a non-target base edge costs one.

\begin{proposition}[lock-placement audit]\label{prop:placements}
Among all \(2^{14}=16{,}384\) subsets \(L\), precisely \(2{,}280\)
make \(\Delta_L(H)>0\) for every extendable \(H\ne T\).
No such set has size at most seven.  Exactly \(180\) have size eight.
For
\[
                         L=\{0,2,3,4,7,8,9,10\},          \tag{18}
\]
the minimum strict gap is one.
\end{proposition}

\begin{proof}[Exact finite proof]
The base cycle-space enumeration gives \(1{,}008\) extendable terminal
supports.  Formula (17) is evaluated for all \(16{,}384\) lock masks
against the other \(1{,}007\) supports.  The asserted counts and gap
follow.  An independent checker enumerates lock sets by combinations
within each cardinality and reproduces the same \(180\) minimum masks.
\end{proof}

\begin{theorem}\label{thm:main}
There is a simple connected bridgeless cubic graph \(G\) on \(162\)
vertices and \(243\) edges whose unique globally minimum extendable
projection has size \(54\) and has no clean extension.  The projection
is the disjoint union of two \(27\)-circuits.
\end{theorem}

\begin{proof}
Use the eight substitutions in (18).  Each pole contributes \(18\)
fresh vertices and replaces one edge by \(28\), so
\[
 |V(G)|=18+8\cdot18=162,\qquad
 |E(G)|=27+8\cdot27=243.                                 \tag{19}
\]
The construction is simple and cubic.  The base and closure are
bridgeless; equivalently, the literal checker deletes every edge of
\(G\) in turn and retains connectivity.

Any nowhere-zero flow on \(G\) contracts through every pole.  Summing
conservation over the pole interior shows that its two nonzero terminal
values agree.  Restoring \(e_*\) gives a nowhere-zero flow on \(S\),
and contracting all poles gives one on \(B\).  Lemma~\ref{lem:lock}
and (17) therefore lower-bound every projection.

The transformed flows (14)--(15) lift (6) and attain
\[
                             14+5\cdot8=54.
\]
Proposition~\ref{prop:placements} forces equality to have terminal
support \(T\).  Equality in each pole then forces the unique local
support (12).  Hence the literal \(54\)-edge projection is unique.
Four locked edges lie on each original \(7\)-circuit, and every
replacement changes one support edge into a \(6\)-edge path.  Each
therefore expands to length \(7+4\cdot5=27\).

It remains to prove uncleanability.  Deleting the \(54\)-edge support
leaves the five old components in (8) and eight pole-interior
components.  Suppose an extension upstairs were clean.  Contract its
poles.  On every old component, a replaced support edge contributes
exactly the colour on its endpoint link; this is the colour of the
contracted edge.  Thus all old component cut-colour parities are
unchanged by contraction.  A clean extension upstairs would therefore
give a clean extension of \(T\) on \(B\), contradicting
Proposition~\ref{prop:base}.
\end{proof}

\section{The graph still has a five-cycle double cover}

For every edge \(e\) and index \(i\in\{0,\ldots,4\}\), introduce
\(x_{e,i}\in\{0,1\}\).  The standard mandatory encoding is
\[
 \sum_{i=0}^4x_{e,i}=2\quad(e\in E),\qquad
 \sum_{e\ni v}x_{e,i}=0\pmod2\quad(v\in V,\ 0\le i<5).    \tag{20}
\]
The \(i\)-th selected edge set is Eulerian, and the edge equations give
exact double coverage.  Conversely every indexed five-cycle double
cover gives such an assignment.  At a cubic vertex the three incident
two-subset labels are the three pairs on some three-element subset of
\(\{0,\ldots,4\}\).

There is a short human construction for \(G\).  In the edge orders
(5) and (10), the following rows are five-cover pair labels:
\[
\begin{array}{ll}
B:&
\texttt{34 23 24 14 13 12 13 03 04 02 01 04 03 02}\\
 &\texttt{14 24 34 12 34 23 24 23 34 23 24 23 34},\\[1mm]
S:&
\texttt{24 23 12 13 13 01 14 34 13 24 14 12 04 34}\\
 &\texttt{03 34 13 14 03 02 23 02 04 24 24 23 34}.
\end{array}                                                \tag{21}
\]
Each pair records the two cover indices containing that edge.  Direct
inspection at each cubic vertex checks that the three incident pairs
form a triangle.

For a replaced base edge with pair \(P\), permute the five indices on
its copy of \(S\) so that the pair on \(e_*\) becomes \(P\).  Delete
the base edge and \(e_*\), and label both new terminal links by \(P\).
At every exposed endpoint, a removed incidence is replaced by the same
pair, so all five parity equations remain even.  Applying this operation
independently to the eight locks proves:

\begin{proposition}\label{prop:positive}
The graph \(G\) of Theorem~\ref{thm:main} has a five-cycle double cover.
\end{proposition}

The artifact package also contains a separately generated literal
assignment for all \(243\) edges.  It is checked directly against (20)
and against a \(1{,}215\)-variable, \(6{,}885\)-clause CNF expansion.
The five class sizes in one independent certificate are
\[
                         (92,109,97,94,94),                \tag{22}
\]
whose sum is \(2|E(G)|=486\).  This is a positive SAT certificate, not
an UNSAT claim.

\section{Verification and trust boundary}

The primary package freezes the labelled \(243\)-edge list, canonical
graph6 and sparse6 encodings, two five-cover certificates, the checker
source, replay output, and a SHA-256 ledger.  A second package rebuilds
the graph in a different edge order.  Its cycle-space code uses
incidence-matrix row reduction instead of fundamental circuits, its
bridge test uses Tarjan low links instead of repeated edge deletion, and
its lock enumeration is organized by combinations instead of integer
masks.  The two constructed graphs canonicalize to the same graph6
record.

\begin{center}
\begin{tabular}{@{}lr@{}}
\toprule
Quantity & Independently reproduced value\\
\midrule
Closure cycles & \(1{,}024\)\\
Closure exact low-zero sets & \(70{,}780\)\\
\((a_0,a_1)\) & \((7,5)\)\\
Base extendable supports & \(1{,}008\)\\
Base extensions of \(T\) & \(15{,}360\)\\
Clean base extensions & \(0\)\\
Lock masks & \(16{,}384\)\\
Minimum eight-lock masks & \(180\)\\
Order and size of \(G\) & \(162,243\)\\
\(\mu(G)\) and minimum-support ties & \(54,1\)\\
Components of \(G-H\) & \(13\)\\
FiveCDC XOR rows & \(810\)\\
\bottomrule
\end{tabular}
\end{center}

The human argument reduces the theorem to these finite counts and
literal certificates.  It does not replace exact enumeration of the two
ten-dimensional cycle spaces or of the \(2^{14}\) lock masks.  The
checkers use only the Python standard library; canonicalization uses
Brendan McKay's \texttt{labelg}.  Full reproduction commands and frozen
digests accompany the source.

\section{Relation to prior work and novelty caution}

The ordinary Cycle Double Cover Conjecture is distinct from the
five-member strengthening.  OpenAI announced a proof of the ordinary
conjecture in July 2026~\cite{OpenAI2026CDC}; expositions by Oum and
Geelen clarify that argument~\cite{Oum2026CDC,Geelen2026CDC}.  Oum
records the five-member version separately.  The present graph has a
five-member cover, so it conflicts with none of those results.

Hu\v{s}ek and \v{S}\'amal give the immediate clean-flow criterion and
state the five-member problem in that language~\cite{HusekSamal2026}.
Hoffmann-Ostenhof gives an earlier prescribed-subgraph
characterization~\cite{HoffmannOstenhof2013}.  Jin, Mazzuoccolo, and
Steffen study other Fano-flow conjectures through cores, joins, and weak
oddness, and notably illustrate that natural strengthening principles in
this area can fail~\cite{JinMazzuoccoloSteffen2018}.  Their objectives
are not the minimum extendable binary projection \(\mu(G)\) used here.

The master lemma is an elementary reparameterization of nowhere-zero
\(\F_2^3\)-flows by an exact zero set and a join; no priority is claimed
for that identity.  The specific strict-lock construction and the
\(162\)-vertex counterexample appear new to this project, but the
literature screen was performed immediately after the Hu\v{s}ek--\v{S}\'amal
preprint appeared.  It cannot establish priority or exclude unpublished
parallel work.  Expert comparison and independent human peer review are
required before any novelty claim.

\section*{AI-use disclosure}

OpenAI Codex agents, directed by Atharva Vaidya, developed the
minimum-projection route, found and minimized the strict two-pole lock,
implemented the enumerations, generated the certificates, and drafted
this manuscript.  Separate Codex agents wrote independent checkers and
performed hostile audits.  Atharva Vaidya directed the project and is
responsible for deciding whether and how to disseminate the result.
The work has not received independent human peer review.  All finite
claims, source code, literal graph encodings, and positive cover
certificates are supplied for human verification.

\bibliographystyle{plain}
\bibliography{references}

\end{document}
